\begin{abstract}
    Modern GPU-centric systems run heterogeneous workloads whose performance and efficiency largely depend on policies for memory placement, kernel scheduling, and observability. Today, these policies are either implemented inside framework-specific runtimes or hard-coded into proprietary vendor stacks, without programmable OS-level interface. Consequently, adapting to changing workloads, SLAs, or hardware still requires invasive code changes and operational disruptions. We argue that the OS GPU driver and device layer, uniquely positioned with cross-tenant visibility and fine-grained control over privileged mechanisms, must serve as a programmable policy narrow waist. However, extending CPU-side programmability frameworks like eBPF to GPUs is challenging: GPU drivers tightly couple resource management with low-level hardware interactions, and host-only or device-only approaches fail to deliver coherent, cross-layer policies.
    
    We present \sys, an eBPF-based policy substrate that spans CPU hosts and GPU devices. \sys refactors GPU drivers to expose safe, verified eBPF hooks at key control points, and allow offloads policies onto GPU devices via a lightweight JIT, coordinated by cross-layer states and single control plane. This unified runtime allows workload-specific memory offloading, kernel preemption, and observability policies to be expressed and enforced consistently across kernel and device contexts, providing strong safety guarantees and low overhead. We implement \sys by extending Linux GPU drivers and building a GPU-side runtime. Evaluation on realistic workloads, including inference, training, and vector search, shows that \sys-based policies improve throughput by up to $5\times$, reduce tail latency by up to $2\times$, and incur only a few percent observability overhead. These results demonstrate the practicality of a driver-anchored, programmable CPU-GPU policy substrate.
\end{abstract}
\section{Introduction}

GPUs now dominate the compute and capital costs of modern computing systems, powering workloads that range from latency-sensitive LLM inference pipelines to large-scale training, graph analytics, and vector search. Across these workloads, end-to-end performance and efficiency depend not just on raw FLOPs, but on \emph{policies}: how the system places and migrates data across device and host memory, admits and schedules kernels in multi-tenant environments, and collects observability signals for online adaptation. These policies must evolve as models, SLAs, and hardware change.

On the CPU side, systems have already transitioned toward programmable policies. The extended Berkeley Packet Filter (eBPF) ecosystem supports dynamically loaded, verified programs for scheduling (\texttt{sched\_ext}), page-cache eviction (\texttt{cache\_ext}), TCP congestion control (\texttt{struct\_ops}), and rich observability~\cite{gregg16,gregg2021computing,Bachl21,Zhong22,jia2023programmable,schedext-docs,cacheext-docs}. These designs expose stable kernel mechanisms via structured attach points, enabling dynamic policy updates without kernel reboots or patches. Treating GPU resource management as another kernel subsystem and attaching host-side eBPF programs to GPU driver events is thus an appealing starting point.

The GPU stack, however, has evolved differently. Vendor drivers and firmware implement fixed, opaque mechanisms for memory management and offloading, context scheduling, and performance event collection. GPU-only instrumentation frameworks such as CUPTI, NVBit, Neutrino, and eGPU provide telemetry via PTX or SASS-level rewriting, but they primarily target profiling, lack unified safety guarantees, and do not integrate with privileged host mechanisms~\cite{cupti,villa2019nvbit,neutrino,egpu}. They cannot control many low-level mechanisms since they happends on CPU, and they cannot enforce online scheduling or memory policies.

More recent serving systems try to retrofit programmability above the driver. Paella exposes software-defined GPU scheduling for model serving~\cite{paella}, Pie lets developers write “inferlets’’ that program the LLM generation loop and KV-cache management~\cite{pie}, and KTransformers offers a flexible CPU–GPU hybrid inference runtime with pluggable scheduling and expert-placement strategies~\cite{ktransformers}. XSched, positioned as an XPU scheduling framework, implements flexible, user-defined policies on top of its XQueue abstraction~\cite{xsched}. In this sense, many user-space GPU systems are already programmable frameworks.

Prior OS-level work moved part of this control into the kernel. TimeGraph and Gdev shifted GPU scheduling and memory control into Linux device drivers~\cite{timegraph,gdev}. More recent systems based on kernel-space interception and preemptive scheduling further refactor driver-level control paths~\cite{kernelintercept,gpreempt}, and LithOS goes even further by re-implementing portions of the userspace GPU runtime stack as an OS component~\cite{lithos}. However, these systems still embed fixed policies in C code and do not expose a general, safe, programmable interface to underlying driver and device mechanisms.

Across these approaches, two fundamental limitations remain for GPU resource management. First, programmability is \emph{framework-scoped}: policies are written against a particular serving system or runtime (e.g., Pie, Paella, XSched, KTransformers) and cannot be reused across arbitrary applications or libraries without porting those applications onto the same stack. Second, user-space systems operate strictly above vendor drivers and firmware and therefore lack a trusted, privileged interface to low-level mechanisms such as page fault handling, MMU updates, interrupt delivery, and device context management. As a result, user-space runtimes are forced to defensively “hoard’’ resources. For example, applications statically pre-allocate GPU memory at startup to manage KV-cache or model weights~\cite{lmcache}, preventing dynamic sharing with other processes. Hardware partitioning techniques like MIG enforce fixed boundaries that cannot adapt to bursty traffic. User-space schedulers are confined to \emph{intra-process} or \emph{intra-framework} decisions; they cannot mitigate interference from “noisy neighbors’’ (other tenants or frameworks) contending for global bandwidth. Even within a single process, they are unprivileged: they cannot intervene at the granularity of hardware interrupts or directly manipulate physical memory mappings, preventing precise compute–transfer overlap and true preemption, and they often impose substantial performance overhead. Finally, deploying new policies typically necessitates invasive changes to the serving framework and service restarts.

Our key observation is that GPU deployments need an OS-level narrow waist for policy. The GPU driver and device execution environment are uniquely positioned with global cross-tenant visibility and direct control over privileged mechanisms such as page migration, kernel scheduling, and preemption. We therefore treat GPU resource management as a first-class kernel subsystem and make it programmable by exposing a small set of verifiable policy interfaces. We further argue that the right abstraction for GPU-centric programmability treats CPU hosts and GPU devices as a single coherent policy domain, supported by a unified intermediate representation, shared state abstraction, and common control plane with lifecycle management. eBPF’s restricted instruction set, verifier guarantees, and rich ecosystem uniquely enable such a unified GPU policy substrate, extending programmable policies from CPUs to GPUs.

However, directly applying host-only eBPF to GPUs fails for two fundamental reasons. First, critical decisions such as thread-block scheduling, warp prioritization, or identifying pathological CUDA kernel behaviors occur device-side within GPU kernels and firmware, invisible to host-only instrumentation. Second, one-size-fits-all GPU drivers tightly couple resource management with hardware interactions (e.g., MMU operations, page fault handling, context management) without stable, safe hooks; naive insertion of probes risks GPU hangs, kernel panics, or resource leaks. A practical GPU policy substrate must therefore separate mechanism from policy inside performance-critical GPU drivers without destabilizing hardware interactions, and must safely execute policy logic within massively parallel, performance-sensitive GPU kernels.

Designing a general substrate for GPU policy programmability introduces four challenges. \emph{C1: Safe and efficient driver extensibility.} We must retrofit mechanism–policy separation into large, performance-critical GPU drivers without destabilizing their low-level interactions with hardware. \emph{C2: Safe and efficient device-side execution.} Policy logic must execute directly within GPU kernels on critical memory access and synchronization paths, while ensuring bounded execution and memory safety in a massively parallel environment. \emph{C3: Cross-layer state and abstractions.} Policies often require correlating host-side and device-side signals (e.g., page access patterns with warp-level stalls), necessitating a unified state abstraction (such as time and virtual address) and efficient synchronization across the CPU–GPU boundary. \emph{C4: A single control plane.} The policy runtime must support unified verification, deployment, and debugging across both host and device execution contexts to avoid introducing additional fragmentation.

We present \sys, a unified, safe, and efficient policy runtime for GPU-centric systems designed to address these challenges. \sys exposes a small, stable set of GPU mechanisms for memory placement and migration, kernel admission and preemption, and kernel-level instrumentation as programmable eBPF attach points within an extended GPU driver. Policies are written as eBPF programs using a restricted instruction subset and are statically verified for memory safety and bounded execution. The same eBPF intermediate representation is then offloaded onto GPU devices: \sys compiles verified bytecode into a GPU-specific device code representation, injects it via lightweight trampolines at kernel launch time, and executes it within a GPU-side runtime. Both host-side and device-side policy programs share a common BPF map abstraction backed by unified memory, enabling low-latency, coherent state sharing without explicit data marshalling. A single verifier and control plane manage policy deployment and lifecycle across host and device, forming a unified policy domain.

We implement \sys on Linux by extending \sysdriver using NVIDIA's open GPU kernel modules~\cite{} and by constructing a GPU-side device JIT runtime. Driver modifications adopt a \texttt{struct\_ops}-style approach, introducing carefully chosen attach points within memory management and scheduling paths and delegating policy decisions to verified eBPF callbacks. On the GPU device, lightweight trampolines save minimal register state, invoke translated eBPF snippets, and restore state, introducing sub-microsecond overhead. Unlike prior GPU observability tools~\cite{cupti,villa2019nvbit,neutrino,egpu} or user-space eBPF implementations, \sys provides a single unified, framework-agnostic policy abstraction across driver and device contexts under one verifier and control plane.

We evaluate \sys across realistic GPU workloads including LLM inference pipelines, GNN training, vector search, and mixed-priority multi-tenant scenarios. Using \sys, we implement adaptive memory prefetching and eviction policies, fine-grained kernel preemption strategies, and bcc-style GPU kernel observability tools. \sys-based policies improve throughput by up to $5\times$ and reduce tail latency by up to $2\times$ compared to static heuristics, while instrumentation-only deployments incur only $2$–$3\%$ overhead. These results confirm that unified, eBPF-based policy programmability spanning CPU and GPU contexts is a practical and powerful abstraction for managing GPU-centric systems.

\begin{itemize}[leftmargin=*,itemsep=0pt]
  \item We identify and characterize key programmability gaps in GPU stacks for memory management, scheduling, and observability, and show how the lack of a programmable, kernel-anchored policy waist at the GPU driver and device layer makes it difficult to adapt to evolving workloads and SLAs.
  \item We design and implement \sys, an eBPF-based policy runtime that treats the GPU driver and device execution layer as a unified OS subsystem. \sys exposes verified driver hooks for memory and scheduling, executes the same eBPF IR inside GPU kernels via a device-side JIT, and shares state through unified-memory-backed BPF maps under a single verifier and control plane.
  \item We demonstrate that \sys enables adaptive, workload-specific policies that significantly improve performance, resource utilization, and tail latency on realistic LLM inference, training, and vector-search workloads, with low overhead and without requiring application or driver restarts.
\end{itemize}